% sec04b-seed.tex — Section 4, expanded seed subsection.
% Writer: 作家_种子 (Writer_Seed). Body LaTeX only; macros live in main.tex.
% This file REPLACES the old short seed subsection. It is \input inside
% sec04-superspace.tex after the selection-rule subsection.

\subsection{The seed channel $(b,b)$: the complete one-loop computation}%
\label{sec:superspace:seed}

We now compute the seed channel $(b,b)$, the simplest nonzero channel and the
scale anchor of the whole table of Section~\ref{sec:superspace:table}, at full
granularity: the graph inventory with orientations, the Wick contraction
enumeration, the amplitude assembly, the $D$-algebra numerator chain, the
evanescent extraction, the pole cancellation, the loop integral, and the
restoration of the letter normalization, closing with the word-by-word
comparison against the holomorphic-twist result of~\cite{BK}. We fix once and
for all an orientation of the two field-strength insertions and a marked
internal line. The graph weights below are computed with unit-normalized
field-strength insertions $\nabla_+\cW_+$; the letter normalization
$b=-(i/\sqrt2)\nabla_+\cW_+$ is restored at the end.

\paragraph{1. Setup: the channel, the letters, and the marking conventions.}
The channel is the one-loop Schwinger--Dyson descent of the ordered composite
of two field-strength insertions,
\begin{equation}\label{eq:ss:seedchannel}
\nabla_{-}\Big(\nabla_{+}\cW_{+}^{A}\,\nabla_{+}\cW_{+}^{B}\Big),
\qquad
b^{A}=-\frac{i}{\sqrt2}\,\nabla_{+}\cW_{+}^{A},
\end{equation}
with the letter $b$ of Definition~\ref{def:ss:letters}. Both insertions are
Grassmann-even, so the Leibniz rule~\eqref{eq:ss:leibniz} contributes no sign
and splits the probe into two marked occurrences,
\begin{equation}\label{eq:ss:seedleibniz}
\nabla_{-}\Big(\nabla_{+}\cW_{+}^{A}\,\nabla_{+}\cW_{+}^{B}\Big)
=\big(\nabla_{-}\nabla_{+}\cW_{+}^{A}\big)\,\nabla_{+}\cW_{+}^{B}
+\nabla_{+}\cW_{+}^{A}\,\big(\nabla_{-}\nabla_{+}\cW_{+}^{B}\big) .
\end{equation}
These are two different marked occurrences, not a multiplicity of either one;
there is no factor $2$ from the source. We fix the ordered external data once
and for all: an antichiral background leg $\bcW^{D}(q)$ at the antichiral
cubic vertex $v_1$, a chiral background leg $\cW_{+}^{E}(p)$ at the chiral
cubic vertex $v_2$, and the composite source vertex $S$ carrying both
insertions and the $\nabla_{-}$ probe. All momenta are incoming
(Remark~\ref{rem:ss:nevercontract}), the loop momentum is $k$, and the three
internal edges are routed as
\begin{equation}\label{eq:ss:seedrouting}
r_0=k\quad(S\to v_1),\qquad
r_1=k+q\quad(v_1\to v_2),\qquad
r_2=k+p+q\quad(v_2\to S).
\end{equation}
There are exactly two ordered source attachments: the \emph{directed}
attachment, in which the insertion of color $A$ contracts toward the
antichiral vertex $v_1$ and the insertion of color $B$ toward the chiral
vertex $v_2$, and the \emph{reflected} attachment, in which they are
interchanged. The complete set of source/mark rows is
\begin{center}
\begin{tabular}{cccl}
\hline
source attachment & marked occurrence & rank-two evanescent numerator &
ordered output\\
\hline
$A\to\bcW,\ B\to\cW$ & $A$ & yes &
$\big\langle\bcW,\nabla_{+}\cW_{+}\big\rangle$\\
$A\to\bcW,\ B\to\cW$ & $B$ & no & none\\
$A\to\cW,\ B\to\bcW$ & $A$ & no & none\\
$A\to\cW,\ B\to\bcW$ & $B$ & yes &
$\big\langle\nabla_{+}\cW_{+},\bcW\big\rangle$\\
\hline
\end{tabular}
\end{center}
The two anomaly-bearing rows are the directed and the reflected term; they
are not two copies of either term. Reflection is a distinct ordered output
with weight one.

\paragraph{2. Graph inventory.}
There is exactly one gauge-sector graph topology at order $g^{2}$: the mixed-chirality
pure-gauge triangle, with one chiral cubic vertex, one antichiral cubic
vertex, the composite source vertex, and three internal vector-superfield
lines. With the canonically normalized prepotential
$u:=\cV/(\sqrt2\,g)$ and $[X,u]^{A}:=f^{ABC}X^{B}u^{C}$, the two cubic
actions generating the vertices are
\begin{equation}\label{eq:ss:seedcubic}
S_{+,3}=-\frac{g}{4}\int d^{8}z\;\cW^{A\alpha}\,[\cD_{\alpha}u,u]^{A},
\qquad
S_{-,3}=+\frac{g}{4}\int d^{8}z\;\bcW^{A}_{\dot\alpha}\,
[\bcD^{\dot\alpha}u,u]^{A}.
\end{equation}
The two ordered chiral assignments of the interaction expansion collapse:
\begin{equation}\label{eq:ss:seedactionorder}
\frac{1}{2!}\left(S_{+,3}S_{-,3}+S_{-,3}S_{+,3}\right)=S_{+,3}S_{-,3},
\end{equation}
so the action-order factor is exactly $1$. The two ordered parent triangles
are displayed in Figure~\ref{fig:ss:seedtriangles}.

\begin{figure}[t]
\centering
\begin{tikzpicture}
  \node[svertex, label=above:$S$] (S) at (90:1.6) {};
  \node[svertex, label=below left:$v_1$] (v1) at (210:1.6) {};
  \node[svertex, label=below right:$v_2$] (v2) at (330:1.6) {};
  \draw[sV] (S) to node[midway,left,align=center,font=\scriptsize]
        {$r_0{=}k$\\$[A]$} (v1);
  \draw[sV] (v1) to node[midway,below,font=\scriptsize] {$r_1{=}k{+}q$} (v2);
  \draw[sV] (v2) to node[midway,right,align=center,font=\scriptsize]
        {$r_2{=}k{+}p{+}q$\\$[B]$} (S);
  \draw[sleg] (v1) -- ++(210:1.15)
        node[below left,font=\scriptsize] {$\bcW^{D}(q)$};
  \draw[sleg] (v2) -- ++(330:1.15)
        node[below right,font=\scriptsize] {$\cW_{+}^{E}(p)$};
  \draw[sleg] (S) -- ++(150:0.95)
        node[sinsert,label={[font=\scriptsize]above left:%
        {$\nabla_{+}\cW_{+}^{A}$ (marked)}}] {};
  \draw[sleg] (S) -- ++(30:0.95)
        node[sinsert,label={[font=\scriptsize]above right:%
        {$\nabla_{+}\cW_{+}^{B}$}}] {};
  \draw[sleg] (S) -- ++(90:0.95) node[above,font=\scriptsize] {$\nabla_{-}$};
  \node[font=\small] at (0,-3.1) {(a) directed: $(A\to\bcW,\ B\to\cW)$,
        probe on $A$};
\end{tikzpicture}
\hspace{1.2em}
\begin{tikzpicture}
  \node[svertex, label=above:$S$] (S) at (90:1.6) {};
  \node[svertex, label=below left:$v_1$] (v1) at (210:1.6) {};
  \node[svertex, label=below right:$v_2$] (v2) at (330:1.6) {};
  \draw[sV] (S) to node[midway,left,align=center,font=\scriptsize]
        {$r_0{=}k$\\$[B]$} (v1);
  \draw[sV] (v1) to node[midway,below,font=\scriptsize] {$r_1{=}k{+}q$} (v2);
  \draw[sV] (v2) to node[midway,right,align=center,font=\scriptsize]
        {$r_2{=}k{+}p{+}q$\\$[A]$} (S);
  \draw[sleg] (v1) -- ++(210:1.15)
        node[below left,font=\scriptsize] {$\bcW^{D}(q)$};
  \draw[sleg] (v2) -- ++(330:1.15)
        node[below right,font=\scriptsize] {$\cW_{+}^{E}(p)$};
  \draw[sleg] (S) -- ++(150:0.95)
        node[sinsert,label={[font=\scriptsize]above left:%
        {$\nabla_{+}\cW_{+}^{A}$}}] {};
  \draw[sleg] (S) -- ++(30:0.95)
        node[sinsert,label={[font=\scriptsize]above right:%
        {$\nabla_{+}\cW_{+}^{B}$ (marked)}}] {};
  \draw[sleg] (S) -- ++(90:0.95) node[above,font=\scriptsize] {$\nabla_{-}$};
  \node[font=\small] at (0,-3.1) {(b) reflected: $(A\to\cW,\ B\to\bcW)$,
        probe on $B$};
\end{tikzpicture}
\caption{The two ordered parent triangles of the seed channel $(b,b)$. Each
triangle has the composite source vertex $S$ carrying the two letter
insertions (filled dots) and the $\nabla_{-}$ probe leg, the antichiral cubic
vertex $v_1$ with external leg $\bcW^{D}(q)$, the chiral cubic vertex $v_2$
with external leg $\cW_{+}^{E}(p)$, and three internal vector-superfield
lines (wavy) with edge momenta $r_0=k$, $r_1=k+q$, $r_2=k+p+q$. The bracket
$[A]$/$[B]$ on an edge records which source color contracts toward that
vertex. (a) The directed orientation, with the probe on the $A$ occurrence.
(b) The reflected orientation, with the probe on the $B$ occurrence.}
\label{fig:ss:seedtriangles}
\end{figure}

For each marked occurrence, the Schwinger cut of
Section~\ref{sec:superspace:cutting} collapses the internal edge carrying the
marked line: the occurrence probed at the $v_1$-side port collapses the edge
$e_0$ with momentum $r_0$, and the occurrence probed at the $v_2$-side port
collapses the edge $e_2$ with momentum $r_2$. For one marked edge the row
count is
\begin{equation}\label{eq:ss:seedcontactcount}
\frac{1}{2!}\cdot 2_{\text{action orderings}}\cdot
2_{\text{Wick endpoints}}=2\ \text{endpoint rows},
\end{equation}
so there are exactly four occurrence-resolved Schwinger contact rows in total
(two cut edges times two Wick endpoints). The contact topology is displayed
in Figure~\ref{fig:ss:seedcontact}.

\begin{figure}[t]
\centering
\begin{tikzpicture}
  \node[svertex, label=above:$v_c$] (vc) at (0,0) {};
  \node[svertex, label=above:$v_2$] (v2) at (3.6,0) {};
  \draw[sV, bend left=30] (vc) to node[midway,above,font=\scriptsize]
        {$r_1{=}k{+}q$} (v2);
  \draw[sV, bend right=30] (vc) to node[midway,below,font=\scriptsize]
        {$r_2{=}k{+}p{+}q$} (v2);
  % sVcut rendering: main.tex's sVcut (snake + markings postaction) triggers a
  % PGF ``Dimension too large'' error under pgf 3.1.10, so the identical cut-line
  % rendering is drawn here as an sV wavy line with an explicit $\times$ marker.
  \draw[sV] (vc) -- ++(195:1.2) coordinate (cutend);
  \node[font=\scriptsize] at ($(vc)!0.5!(cutend)$) {$\times$};
  \node[below left,font=\scriptsize] at (cutend) {cut $e_0$};
  \draw[sleg] (vc) -- ++(150:1.05)
        node[sinsert,label={[font=\scriptsize]above left:%
        {$\nabla_{+}\cW_{+}^{B}$}}] {};
  \draw[sleg] (vc) -- ++(120:1.05)
        node[above,font=\scriptsize] {$\bcW^{D}(q)$};
  \draw[sleg] (v2) -- ++(0:1.2)
        node[right,font=\scriptsize] {$\cW_{+}^{E}(p)$};
\end{tikzpicture}
\caption{The Schwinger contact (cut) topology of the seed channel. Cutting
the marked edge $e_0$ collapses the source vertex $S$ and the antichiral
vertex $v_1$ into the single contact vertex $v_c$: the probed occurrence
$\nabla_{-}(\nabla_{+}\cW_{+}^{A})$ has become the Euler-operator insertion
that collapses the $e_0$ line (wavy with $\times$), while the unmarked
insertion $\nabla_{+}\cW_{+}^{B}$ and the external leg $\bcW^{D}(q)$ remain
at $v_c$. Two internal vector lines with momenta $r_1=k+q$ and $r_2=k+p+q$
run from $v_c$ to the chiral vertex $v_2$. The $e_2$ family is the mirror
image, with contact vertex $v_c=S\sqcup v_2$ and internal lines $r_0$ and
$r_1$. With the two Wick port rows per edge this gives the four
occurrence-resolved contact rows of paragraph~7.}
\label{fig:ss:seedcontact}
\end{figure}

Every remaining order-$g^{2}$ orbit vanishes exactly in the anomaly sector.
The spectator self-energy insertions, the one-cubic-vertex bubbles, and the
quartic seagull obey
\begin{align}
\big\langle I_2\big\rangle_0 &= 0
&&\text{(source-quadratic self-energy: its only denominator is $k^{2}$,
hence a scaleless tadpole),}\notag\\
-\frac{1}{\h}\big\langle I_1S_{g3}\big\rangle_0 &= 0
&&\text{(one-cubic bubble: numerator trace }\ \operatorname{tr}_4 H=2-2=0
\text{),}\label{eq:ss:seedorbits}\\
-\frac{1}{\h}\big\langle I_0S_{g4}\big\rangle_0 &= 0
&&\text{(quartic seagull: }\ \operatorname{tr}_4 H_4=-4+4=0\text{),}\notag
\end{align}
so only the two-cubic triangle
$(2\h^{2})^{-1}\langle I_0S_{g3}S_{g3}\rangle_0$ of
Figure~\ref{fig:ss:seedtriangles} and its Schwinger cuts survive. The matter
orbits $I_2$, $I_1S_{m3}$, $I_0S_{m4}$ vanish in the same way; the matter
triangle of the same channel is evaluated in paragraph~9.

\paragraph{3. Wick contractions.}
The three internal lines are vector-superfield contractions with propagator
\begin{equation}\label{eq:ss:seeduprop}
\big\langle u^{A}(p,1)\,u^{B}(-p,2)\big\rangle
=-\frac{\h\,\delta^{AB}}{p^{2}}\,\delta^{4}(\theta_1-\theta_2),
\end{equation}
with the Killing form normalized to $\kappa_{AB}=\delta_{AB}$. The two cubic
vertices each supply two quantum-port assignments through the derivative
differences of their second variations, the Hessians
\begin{equation}\label{eq:ss:seedhess}
H_{\cW}=+\frac{ig}{4}\,f^{UCE}\,\cW^{E\gamma}
\big(\cD_{C\gamma}-\cD_{U\gamma}\big),
\qquad
H_{\bcW}=-\frac{ig}{4}\,f^{UCD}\,\bcW^{D}_{\dot\gamma}
\big(\bcD_{C}{}^{\dot\gamma}-\bcD_{U}{}^{\dot\gamma}\big).
\end{equation}
The labeled double variation is
\begin{equation}\label{eq:ss:seedports}
\delta_1\delta_2[\cD u,u]=[\cD u_1,u_2]+[\cD u_2,u_1],
\end{equation}
whose two summands are the port rows $(S,B)$ and $(B,S)$ --- the row in which
the first, respectively the second, quantum port carries the source role.
The second variation already contains the two labeled quantum-port
assignments, so no $2!$ from the equal Hessians remains. For one fixed
directed source attachment the port-preserving Wick factor is likewise $1$:
the two rows are equal by relabeling of the Gaussian integral and sum against
the $1/2!$ of the second variation, giving the Wick weight
\begin{equation}\label{eq:ss:wick}
w_{\mathrm{Wick}}=\frac{1}{2!}(1+1)=1 .
\end{equation}

\begin{remark}[Granularity of the contraction enumeration]%
\label{rem:ss:seedwick}
The enumeration above is organized by the quantum-role rows $(S,B),(B,S)$
and by the derivative endpoint pairs $(r_i,r_j)$; a field-by-field pairing
picture of the individual $u$ fields is not used at any step. The endpoint
content is not multiplied after the two Hessians are expanded: no endpoint
row acquires an additional factor.
\end{remark}

\paragraph{4. Amplitude assembly.}
The linearized unit insertion and its probed descendant are
\begin{equation}\label{eq:ss:seedinsert}
K_{+}^{c}:=-\frac{1}{4\sqrt2}\,\cD_{+}\bcD^{2}\cD_{+},
\qquad
\nabla_{-}K_{+}^{c}
=-\frac{1}{4\sqrt2}\,\cD_{-}\cD_{+}\bcD^{2}\cD_{+}
=-\frac{1}{8\sqrt2}\,\cD^{2}\bcD^{2}\cD_{+},
\end{equation}
where on the internal lines the frame derivatives act as their flat
counterparts $\cD_\pm$, $\bcD_{\dot\pm}$ and we used the spinor-derivative
square move
\begin{equation}\label{eq:ss:seedDmoves}
\cD^{2}=2\,\cD_{-}\cD_{+},
\qquad
\bcD^{2}=2\,\bcD_{\dot+}\bcD_{\dot-}.
\end{equation}
After the three contractions, the unmarked leg retains the residual operator
$K_{+}^{c}$ and the probed leg retains $\nabla_{-}K_{+}^{c}$; their numerical
product is
\begin{equation}\label{eq:ss:seedinsertfactor}
\left(-\frac{1}{8\sqrt2}\right)\left(-\frac{1}{4\sqrt2}\right)=\frac{1}{64}.
\end{equation}
The vertex (exponent) Hessians of~\eqref{eq:ss:seedhess}, including the
$1/\h$ of each interaction vertex, are
\begin{equation}\label{eq:ss:seedvhess}
\widehat H_{\cW}=-\frac{ig}{4\h}\,f^{UCE}\,\cW^{E\gamma}
\big(\cD_{C\gamma}-\cD_{U\gamma}\big),
\qquad
\widehat H_{\bcW}=+\frac{ig}{4\h}\,f^{UCD}\,\bcW^{D}_{\dot\gamma}
\big(\bcD_{C}{}^{\dot\gamma}-\bcD_{U}{}^{\dot\gamma}\big),
\end{equation}
and assembling them with the three propagators~\eqref{eq:ss:seeduprop} gives
\begin{equation}\label{eq:ss:seedassembly}
\left(-\frac{ig}{4\h}\right)\left(+\frac{ig}{4\h}\right)(-\h)^{3}
=\frac{g^{2}}{16\h^{2}}\,(-\h^{3})
=-\frac{\h g^{2}}{16}.
\end{equation}
The three Killing deltas of the propagators reduce the color factor of the
two Hessians, with the source colors attached per the fixed orientation, to
\begin{equation}\label{eq:ss:seedcolor}
f^{UCD}f^{U'C'E}\,\delta_{UU'}\delta_{CC'}
\;\longrightarrow\;
f^{ACD}f^{BCE}
\qquad(\text{sum over }C).
\end{equation}

\begin{remark}[Two equivalent normalization frames]\label{rem:ss:seedframes}
Equation~\eqref{eq:ss:seedassembly} is the canonical-letter bookkeeping:
vertex magnitudes $g/(4\h)$, propagators $-\h$, and the explicit insertion
factor $1/64$ of~\eqref{eq:ss:seedinsertfactor}. In the unit-insertion frame
used for the remaining steps the same graph weight is organized with the
ordered cubic vertex factors
\begin{equation}\label{eq:ss:vertices}
\left(+\frac{ig}{2}\right)\left(-\frac{ig}{2}\right)=\frac{g^{2}}{4}
\end{equation}
and the $D$-algebra weight $w_D$ of~\eqref{eq:ss:dweight} below. The two
bookkeepings describe the same graph weight: in the canonical frame the
vertex factor sits inside $-\h g^{2}/16$, in the unit-insertion frame inside
$g^{2}/4$, the difference being the propagator rescaling
$u=\cV/(\sqrt2\,g)$ on the three internal lines together with the insertion
rescaling. Each cubic vertex carries $\h^{-1}$ and each of the three
propagators carries $\h$, leaving the net one-loop power $\h$.
\end{remark}

\paragraph{5. The $D$-algebra numerator chain.}
For the fixed orientation the covariant spinor numerator is the chain
\begin{equation}\label{eq:ss:seedspinchain}
N_{G,+}{}^{\dot\alpha}
=(r_0+r_1)_{+\dot\beta}\,(ip)^{\dot\beta\gamma}\,
(r_1+r_2)_{\gamma}{}^{\dot\alpha},
\qquad
T_{m\rho n,+}{}^{\dot\alpha}
:=(\sigma_m\bar\sigma_\rho\sigma_n)_{+}{}^{\dot\alpha},
\end{equation}
where the explicit $i$ is tied to the locked dotted-index convention
$k_{+}^{\dot0}=\sigma_E(k)_{0\dot1}$,
$k_{+}^{\dot1}=-\sigma_E(k)_{0\dot0}$; in the opposite convention the same
covariant chain is written with $p^{\dot\beta\gamma}$ and no explicit $i$.
The $D^{2}/\bcD^{2}$ moves of~\eqref{eq:ss:seedDmoves} put the marked
insertion into the form $-\frac{1}{8\sqrt2}\cD^{2}\bcD^{2}\cD_{+}$ displayed
in~\eqref{eq:ss:seedinsert}. The loop is saturated by the identity
\begin{equation}\label{eq:ss:seedsat}
\delta^{4}(\theta_{12})\,\cD^{2}\bcD^{2}\,\delta^{4}(\theta_{12})
=16\,\delta^{4}(\theta_{12})
\end{equation}
of Section~\ref{sec:gaugebv}, and the two mixed anticommutators
\begin{equation}\label{eq:ss:seedmixed}
\{\cD,\bcD\}=2\,\sigma\!\cdot\!p
\end{equation}
contribute one factor $2$ each, so the $D$-algebra content of the closed
loop is
\begin{equation}\label{eq:ss:seeddalg}
\left(\frac{1}{64}\right)(16)(2)(2)=1 .
\end{equation}
This is a $D$-algebra factor; it is not four additional Wick contractions.
Equivalently, in the unit-insertion frame the closed-loop mixed $D$-algebra
weight is
\begin{equation}\label{eq:ss:dweight}
w_D=\frac{1}{32}\cdot16\cdot2\cdot2=2 ,
\end{equation}
where the $1/32$ is the insertion normalization of that frame and the $16$
is the loop-saturation identity~\eqref{eq:ss:seedsat}. The total prefactor
before the loop integral is therefore
\begin{equation}\label{eq:ss:prefactor}
w_{\mathrm{Wick}}\,\frac{g^{2}}{4}\,w_D
=1\cdot\frac{g^{2}}{4}\cdot2=\frac{g^{2}}{2} .
\end{equation}
The four derivative endpoint rows of one cut word, with their raw and
transport signs, are
\begin{align}
(r_0,r_1):&\quad (-1)(-1)=+1,\notag\\
(r_0,r_2):&\quad (+1)(+1)=+1,\label{eq:ss:seedsigntable}\\
(r_1,r_1):&\quad (+1)(+1)=+1,\notag\\
(r_1,r_2):&\quad (-1)(-1)=+1,\notag
\end{align}
and the closed word carries in addition the single odd source-loop
integration-by-parts sign $s_{\mathrm{global}}=-1$, so the signed row weight
is
\begin{equation}\label{eq:ss:seedrowsign}
w_{T,\mathrm{rows}}
=s_{\mathrm{global}}\left(\frac{1}{64}\right)(16)(2)(2)=(-1)(1)=-1,
\qquad
\left(-\frac{\h g^{2}}{16}\right)w_{T,\mathrm{rows}}=+\frac{\h g^{2}}{16}.
\end{equation}
The rank-two finite spinor word is reduced by the four-dimensional identity
\begin{equation}\label{eq:ss:seedfierz}
\sigma\!\cdot\!L\,\bar\sigma\!\cdot\!p\,\sigma\!\cdot\!L
=2(L\!\cdot\!p)\,\sigma\!\cdot\!L-\bar L^{2}\,\sigma\!\cdot\!p ,
\end{equation}
whose two factors of $2L$ give
$4\,\sigma\!\cdot\!L\,\bar\sigma\!\cdot\!p\,\sigma\!\cdot\!L
=8(L\!\cdot\!p)\,\sigma\!\cdot\!L-4\bar L^{2}\,\sigma\!\cdot\!p$; the first
term belongs to the full-$d$ Schwinger-exact tensor reduction, and the
selected inverse-square parent and its cut differ by
\begin{equation}\label{eq:ss:seedsigmafour}
-4\big(\bar L^{2}-L_d^{2}\big)\,\sigma\!\cdot\!p
=-4\,\mul^{2}\,\sigma\!\cdot\!p ,
\end{equation}
which vanishes identically if $\bar L^{2}$ is replaced by $L_d^{2}$ before
the cut subtraction. For the reflected ordering the dotted-index reordering
gives
\begin{equation}\label{eq:ss:seedreflect}
\big(\partial^{\dot\alpha}\,\nabla_{+}\cW_{+}^{D}\big)\,\bcW_{\dot\alpha}^{E}
=-\,\big(\partial_{\dot\alpha}\,\nabla_{+}\cW_{+}^{D}\big)\,
\bcW^{E\dot\alpha},
\end{equation}
one minus sign per reordering, which is the origin of the antisymmetric
ordered pair of outputs. The exact marked words behind the endpoint table
are
\begin{equation}\label{eq:ss:seedGwords}
\mathcal G_1=(b_{0+}+b_{1+})\,W_{02},
\qquad
\mathcal G_2=(b_{0+}-b_{1+})\,W_{02},
\qquad
W_{02}=r_{0+}\wedge r_{2+},
\end{equation}
with $b_{0+}-b_{1+}=q_{+\dot-}$ purely external, so $\mathcal G_2$ has no
rank-two loop component (the rows marked ``no'' in the table of
paragraph~1), while $b_{0+}+b_{1+}=2L_{+\dot-}+\text{external}$ is rank two
(the rows marked ``yes'').

\begin{remark}[Status of the sign table]\label{rem:ss:seedsigns}
The endpoint-row assignment~\eqref{eq:ss:seedsigntable} and the global sign
$s_{\mathrm{global}}=-1$ are recorded primitives of the computation: no
independent human-readable noncommutative $D$-algebra trace emitting them is
quoted here. They are superseded in effect by the exact exterior-algebra
words~\eqref{eq:ss:seedGwords} and by the machine replays of
Remark~\ref{rem:ss:seedchecks}, which confirm the full factor chain with
zero failures.
\end{remark}

\paragraph{6. Evanescent extraction and tensor reduction.}
On every marked line $e$ of every one-loop graph the cutting-failure
identity~\eqref{eq:ss:cuttingfailure} isolates one universal insertion:
\begin{equation}\label{eq:ss:seededge}
\frac{\bar r_e^{\,2}}{D_0D_1D_2}-\frac{1}{\prod_{j\ne e}D_j}
=\frac{\mul^{2}}{D_0D_1D_2},
\end{equation}
with $D_i=r_{i,d}^{2}$; replacing $\bar r_e^{\,2}$ by $r_{e,d}^{\,2}$ before
the subtraction makes the expression identically zero and removes the
anomaly. With $D_0=\ell^{2}$, $D_1=(\ell+p)^{2}$, $D_2=(\ell+p+q)^{2}$, the
ordered Feynman parametrization is
\begin{equation}\label{eq:ss:feynman}
\frac{1}{D_0D_1D_2}
=2\!\int_{x,y,z\geq0}\!dx\,dy\,dz\;\delta(1-x-y-z)\,\frac{1}{(r^{2}+\Delta)^{3}},
\qquad
\int_{x,y,z\geq0}\!dx\,dy\,dz\;\delta(1-x-y-z)=\frac12 ,
\end{equation}
with the shift, writing $Q=p+q$,
\begin{equation}\label{eq:ss:seedshift}
\ell=k+yq+zQ,
\qquad
\Delta=xyq^{2}+xzQ^{2}+yzp^{2}.
\end{equation}
The two numerator legs $L_1=2k+q$ and $L_2=2k+p+2q$ become
\begin{equation}\label{eq:ss:seedL12}
L_1=2\ell+(1-2y)q-2zQ,
\qquad
L_2=2\ell+(1-2y)q+(1-2z)Q,
\end{equation}
so the odd-$\ell$ terms integrate to zero and only the $\ell_m\ell_n$ piece
survives. Hatted rotational reduction and~\eqref{eq:ss:Jnidentities} give
\begin{equation}\label{eq:ss:tensorreduction}
\int_r\frac{r_mr_n}{(r^{2}+\Delta)^{3}}
=\frac{\widehat\delta_{mn}}{d}\,(J_2-\Delta J_3)
=\frac{\widehat\delta_{mn}}{4}\,J_2 ,
\end{equation}
the $\tfrac14$ being exact since
$(1-\tfrac{\epsilon}{2})/(4-2\epsilon)=\tfrac14$ identically in $\epsilon$.
Consequently the pole part of the parent triangle with numerator legs
$L_1^{m}L_2^{n}$ is
\begin{equation}\label{eq:ss:pole}
\operatorname*{Pole}_{\epsilon=0}
\int_{\ell}\frac{L_1^{m}L_2^{n}}{D_0D_1D_2}
=\frac{\widehat\delta^{mn}}{16\pi^{2}\epsilon} .
\end{equation}
Assembling the prefactor~\eqref{eq:ss:prefactor}, the simplex
factor~\eqref{eq:ss:feynman}, the residue~\eqref{eq:ss:pole}, and the spinor
chain $\sigma_m\bar\sigma_\rho\sigma_n$ contracted with the hatted external
momentum $p^{\rho}$, the parent triangle pole is
\begin{equation}\label{eq:ss:parentpole}
\Gamma_{T}
=\frac{\h g^{2}}{32\pi^{2}\epsilon}\,
f^{ACD}f^{BCE}\,\widehat\delta^{mn}\,
\sigma_m\bar\sigma_\rho\sigma_n\,p^{\rho} .
\end{equation}

\paragraph{7. Schwinger contact rows.}
The four occurrence-resolved contact rows of
Figure~\ref{fig:ss:seedcontact} are
\begin{center}
\begin{tabular}{ccll}
\hline
row & cut edge & marked occurrence & Wick port row\\
\hline
$j=1$ & $e_0$ & $\nabla_{-}(\nabla_{+}\cW_{+}^{A})$ & $(S,B)$\\
$j=2$ & $e_0$ & $\nabla_{-}(\nabla_{+}\cW_{+}^{A})$ & $(B,S)$\\
$j=3$ & $e_2$ & $\nabla_{+}\cW_{+}^{A}\,\nabla_{-}(\nabla_{+}\cW_{+}^{B})$
& $(S,B)$\\
$j=4$ & $e_2$ & $\nabla_{+}\cW_{+}^{A}\,\nabla_{-}(\nabla_{+}\cW_{+}^{B})$
& $(B,S)$\\
\hline
\end{tabular}
\end{center}
and each row carries the same sign chain
\begin{equation}\label{eq:ss:signchain}
\left(-\frac14\right)_{\mathrm{cut}}
(-1)_{\mathrm{endpoint}}
\left(\frac12\right)_{\nabla_{-}\nabla_{+}=\nabla^{2}/2}
(2)_{\mathrm{mixed}}
(2)_{\mathrm{Wick}}
\left(-\frac12\right)_{p_{+}}
=-\frac14 ,
\end{equation}
whose factors are, in order: the cut/insertion coefficient, the endpoint
transport sign, the spinor-derivative square
move~\eqref{eq:ss:seedDmoves}, the mixed anticommutator
factor~\eqref{eq:ss:seedmixed}, the two Wick endpoints, and the plus-frame
momentum projection. There is no extra edge factor $2$ and no metric-trace
factor $d$ or $4$ on the contact side. A single row gives
\begin{equation}\label{eq:ss:singlerow}
\Gamma_{C,j}
=-\frac14\cdot\frac{\h g^{2}}{2}\cdot\frac{1}{16\pi^{2}\epsilon}\,
f^{ACD}f^{BCE}\,\delta_4^{mn}\,
\sigma_m\bar\sigma_\rho\sigma_n\,p^{\rho}
=-\frac{\h g^{2}}{128\pi^{2}\epsilon}\,
f^{ACD}f^{BCE}\,\delta_4^{mn}\,
\sigma_m\bar\sigma_\rho\sigma_n\,p^{\rho} ,
\end{equation}
and the four occurrence-resolved rows sum to
\begin{equation}\label{eq:ss:contactrows}
\Gamma_{C}
=\sum_{j=1}^{4}\Gamma_{C,j}
=-\frac{\h g^{2}}{32\pi^{2}\epsilon}\,
f^{ACD}f^{BCE}\,\delta_4^{mn}\,
\sigma_m\bar\sigma_\rho\sigma_n\,p^{\rho} .
\end{equation}

\paragraph{8. Pole cancellation and the $-2\epsilon$ trace.}
The four-dimensional metric parts of~\eqref{eq:ss:parentpole}
and~\eqref{eq:ss:contactrows} cancel exactly: since
$\widehat\delta^{mn}-\delta_4^{mn}=-\breve\delta^{mn}$,
\begin{equation}\label{eq:ss:polecancel}
\Gamma_{T}+\Gamma_{C}
=-\frac{\h g^{2}}{32\pi^{2}\epsilon}\,
f^{ACD}f^{BCE}\,\breve\delta^{mn}\,
\sigma_m\bar\sigma_\rho\sigma_n\,p^{\rho} ,
\end{equation}
a pure evanescent numerator, precisely as in the bookkeeping principle of
Remark~\ref{rem:ss:polebookkeeping}. The external momentum is hatted, so
contracting with $p^{\rho}$ first and using $\breve\delta\,p=0$ and
$\breve\delta^{mn}\bar\sigma_m\sigma_n=\breve\delta^{m}{}_{m}=2\epsilon$
with the four-dimensional spin algebra~\eqref{eq:ss:sigmaalgebra},
\begin{equation}\label{eq:ss:sigmatrace}
p^{\rho}\,\breve\delta^{mn}\,\sigma_m\bar\sigma_\rho\sigma_n
=2p^{\rho}\breve\delta_{\rho}{}^{n}\sigma_n
-p^{\rho}\sigma_\rho\,\breve\delta^{mn}\bar\sigma_m\sigma_n
=0-(\breve\delta^{m}{}_{m})\,p^{\rho}\sigma_\rho
=-2\epsilon\,p^{\rho}\sigma_\rho .
\end{equation}

\paragraph{9. The loop integral and the complete unit-insertion result.}
The shifted coefficient is independent of the simplex parameters, so
$2\int_{\Sigma_2}1=2(\tfrac12)=1$ by~\eqref{eq:ss:feynman}. Substituting
\eqref{eq:ss:sigmatrace} into~\eqref{eq:ss:polecancel} and using the
evanescent master integral~\eqref{eq:ss:masterintegral},
\begin{equation}\label{eq:ss:seedfinal}
\Gamma_{T}+\Gamma_{C}
=-\frac{\h g^{2}}{32\pi^{2}\epsilon}\,(-2\epsilon)\,
f^{ACD}f^{BCE}\,\sigma_\rho p^{\rho}
=\frac{\h g^{2}}{16\pi^{2}}\,
f^{ACD}f^{BCE}\,\sigma\!\cdot\!p .
\end{equation}
The $\sigma$-chain trace forces the position-space output to be a dotted
bilinear of the form~\eqref{eq:ss:bilinear}: this is the output structure
quoted as grading (iv) in Section~\ref{sec:superspace:selection}. The
directed triangle emits the ordered word
$\langle\bcW^{D},\nabla_{+}\cW_{+}^{E}\rangle$; the reflected triangle emits
the reordered word with the minus sign of~\eqref{eq:ss:seedreflect}; the two
words in which the holomorphic momentum acts along the $q$-leg are total
derivatives (equation-of-motion carriers) and drop out of the physical
output. The gauge sector of the channel is therefore
\begin{equation}\label{eq:ss:seedgauge}
\Gamma^{(1)}_{(b,b),\mathrm{gauge}}
=\frac{\h g^{2}}{16\pi^{2}}\,f^{ACD}f^{BCE}\,
\big\{\big\langle\bcW,\nabla_{+}\cW_{+}\big\rangle
-\big\langle\nabla_{+}\cW_{+},\bcW\big\rangle\big\}^{DE} .
\end{equation}
The matter sector of the same channel is the two-$\widetilde\Phi\cV\Phi$
marked triangle. Its two marked placements give, with the same master
integral~\eqref{eq:ss:masterintegral},
\begin{align}
4\h g^{2}\left(\frac23\right)\left(\frac{1}{32\pi^{2}}\right)
&=\frac{\h g^{2}}{12\pi^{2}},
\label{eq:ss:seedmatterone}\\
4\h g^{2}\left(-\frac16\right)\left(\frac{1}{32\pi^{2}}\right)
&=-\frac{\h g^{2}}{48\pi^{2}},
\label{eq:ss:seedmattertwo}
\end{align}
and the two placements sum to
\begin{equation}\label{eq:ss:seedmatternet}
\frac{\h g^{2}}{12\pi^{2}}-\frac{\h g^{2}}{48\pi^{2}}
=\frac{3\,\h g^{2}}{48\pi^{2}}
=\frac{\h g^{2}}{16\pi^{2}},
\end{equation}
carried by the matter words
$\sum_r\big(\langle\nabla_{+}\Phi_r,\widetilde\Phi^{r}\rangle
-\langle\widetilde\Phi^{r},\nabla_{+}\Phi_r\rangle\big)$. Collecting the
gauge sector~\eqref{eq:ss:seedgauge} and the matter
sector~\eqref{eq:ss:seedmatternet}, the complete channel in unit insertions
is
\begin{equation}\label{eq:ss:seedunitresult}
\begin{aligned}
\nabla_{-}\Big(\nabla_{+}\cW_{+}^{A}\,\nabla_{+}\cW_{+}^{B}\Big)
\Big|_{\text{1-loop}}
=\frac{\h g^{2}}{16\pi^{2}}\,f^{ACD}f^{BCE}\,
\Big\{&\big\langle\bcW,\nabla_{+}\cW_{+}\big\rangle
-\big\langle\nabla_{+}\cW_{+},\bcW\big\rangle\\
+&\sum_{r=1}^{3}\Big(
\big\langle\nabla_{+}\Phi_r,\widetilde\Phi^{r}\big\rangle
-\big\langle\widetilde\Phi^{r},\nabla_{+}\Phi_r\big\rangle\Big)
\Big\}^{DE} ,
\end{aligned}
\end{equation}
with the dotted contraction of~\eqref{eq:ss:bilinear} understood in each
bilinear.

\begin{remark}[The matter source-orientation sign]\label{rem:ss:seedmatter}
The relative size of the two matter placements in~\eqref{eq:ss:seedmatterone}
and~\eqref{eq:ss:seedmattertwo} is exact; the overall matter
source-orientation sign is fixed by the target-blind primitive-sign chain
($\eta_{\mathrm{src}}=+1$), giving the word order displayed
in~\eqref{eq:ss:seedunitresult}. The same sign is consistent with the
accepted component bottom projection of Section~\ref{sec:component} (an exact
projection of the superfield result, not an independent derivation) and is
checked independently by the holomorphic-twist comparison of paragraph~11.
\end{remark}

\paragraph{10. Restoring the letter normalization: the boxed result.}
The letters of Definition~\ref{def:ss:letters} give the substitutions
\begin{align}
\nabla_{+}\cW_{+}&=i\sqrt2\,b,
&\big(\nabla_{+}\cW_{+}^{A}\big)\big(\nabla_{+}\cW_{+}^{B}\big)
&=(i\sqrt2)^{2}\,b^{A}b^{B}=-2\,b^{A}b^{B},
\label{eq:ss:seedlettersub}\\
\big\langle\bcW,\nabla_{+}\cW_{+}\big\rangle
&=\big\langle-i\,\partial c,\,i\sqrt2\,b\big\rangle
=\sqrt2\,\big\langle\partial c,b\big\rangle,
&\big\langle\nabla_{+}\Phi_r,\widetilde\Phi^{r}\big\rangle
&=\big\langle\sqrt2\,\beta_r,\gamma^{r}\big\rangle
=\sqrt2\,\big\langle\beta_r,\gamma^{r}\big\rangle ,
\label{eq:ss:seedwordsub}
\end{align}
so that dividing~\eqref{eq:ss:seedunitresult} by the insertion factor $-2$
yields
\begin{equation}\label{eq:ss:seednetletter}
\nabla_{-}\big(b^{A}b^{B}\big)
=\frac{1}{-2}\cdot\sqrt2\cdot\frac{\h g^{2}}{16\pi^{2}}\,
f^{ACD}f^{BCE}\,\{\cdots\}
=-\frac{\h g^{2}}{16\sqrt2\,\pi^{2}}\,f^{ACD}f^{BCE}\,\{\cdots\} .
\end{equation}
Equivalently, in the bookkeeping of the unit-insertion frame: the two
insertions contribute a factor $(-i/\sqrt2)^{2}=-\tfrac12$ relative to the
unit field-strength normalization used above, while the output bilinears of
the seed channel contribute $\sqrt2$ in the alphabet of
Definition~\ref{def:ss:letters}, giving the net prefactor
$-\tfrac12\cdot\sqrt2\cdot\tfrac{\h g^2}{16\pi^2}
=-\tfrac{\h g^{2}}{16\sqrt2\,\pi^{2}}$. The position-space identity reads
\begin{equation}\label{eq:ss:bbresult}
\boxed{\;
\nabla_{-}\big(b^{A}b^{B}\big)
=-\frac{\h g^{2}}{16\sqrt2\,\pi^{2}}\,
f^{ACD}f^{BCE}\,
\Big\{
\langle\partial c,b\rangle-\langle b,\partial c\rangle
+\sum_{r=1}^{3}\big(
\langle\beta_{r},\gamma^{r}\rangle-\langle\gamma^{r},\beta_{r}\rangle
\big)\Big\}^{DE} .\;}
\end{equation}

\begin{remark}[Independent machine cross-checks]\label{rem:ss:seedchecks}
The factor chain of this subsection was replayed externally on $269$ external
slot assignments with zero failures, and on the full color mask ($9216$
entries, with a sparse replay of $2048$) with zero equality failures.
\end{remark}

\paragraph{11. The holomorphic-twist comparison for this channel.}
Budzik and Kulp compute the same channel in the holomorphic twist; their
zero-derivative display \cite[eq.~(3.64)]{BK} is
\begin{equation}\label{eq:ss:seedbk}
Q_1(b^A b^B)
=\kappa^{2} f_{ACD}f_{BCE} \left[
\partial_{\dot\alpha} c^D \partial^{\dot\alpha} b^E
-\partial_{\dot\alpha} b^D\partial^{\dot\alpha} c^E
+\partial_{\dot\alpha}(\beta_I)^D\partial^{\dot\alpha}(\gamma^I)^E
-\partial_{\dot\alpha} (\gamma^I)^D\partial^{\dot\alpha}(\beta_I)^E
\right],
\end{equation}
with $\kappa$ the dual Coxeter normalization $\Tr t_At_B=\kappa\delta_{AB}$.
The dictionary from the twist letters to the unit-normalized composites of
this section is
\begin{equation}\label{eq:ss:seeddictionary}
b\mapsto-\frac{i\rho}{\sqrt2}\,\nabla_{+}\cW_{+},
\qquad
\beta_r\mapsto\frac{\rho}{\sqrt2}\,\nabla_{+}\Phi_r,
\qquad
\gamma^r\mapsto\rho\,\widetilde\Phi^{r},
\qquad
\partial_{\dot\alpha}c\mapsto i\rho\,\bcW_{\dot\alpha},
\end{equation}
with color, flavor, and dotted-spin frames $R_A{}^P$, $U_I{}^r$,
$S_{\dot\alpha}{}^{\dot\beta}$ and a scale $\rho\neq0$, all set to the
identity after tracking. Applied word by word, the dictionary gives
\begin{align}
\partial_{\dot\alpha} c^D\partial^{\dot\alpha} b^E
&\mapsto
(i\rho)\Big(-\frac{i\rho}{\sqrt2}\Big)\,
\big\langle\bcW^{D},\nabla_{+}\cW_{+}^{E}\big\rangle
=\frac{\rho^{2}}{\sqrt2}\,
\big\langle\bcW^{D},\nabla_{+}\cW_{+}^{E}\big\rangle,
\notag\\
-\,\partial_{\dot\alpha} b^D\partial^{\dot\alpha} c^E
&\mapsto
-\frac{\rho^{2}}{\sqrt2}\,
\big\langle\nabla_{+}\cW_{+}^{D},\bcW^{E}\big\rangle,
\label{eq:ss:seedwordmap}\\
\partial_{\dot\alpha}(\beta_I)^D\partial^{\dot\alpha}(\gamma^I)^E
&\mapsto
\frac{\rho^{2}}{\sqrt2}\,U_I{}^r(U^{-1})_s{}^I\,
\big\langle\nabla_{+}\Phi_r^{D},\widetilde\Phi^{sE}\big\rangle
=\frac{\rho^{2}}{\sqrt2}\sum_r
\big\langle\nabla_{+}\Phi_r^{D},\widetilde\Phi^{rE}\big\rangle,
\notag\\
-\,\partial_{\dot\alpha}(\gamma^I)^D\partial^{\dot\alpha}(\beta_I)^E
&\mapsto
-\frac{\rho^{2}}{\sqrt2}\sum_r
\big\langle\widetilde\Phi^{rD},\nabla_{+}\Phi_r^{E}\big\rangle .
\notag
\end{align}
The input pair contributes $(-i\rho/\sqrt2)^{2}=-\rho^{2}/2$, undoing it
contributes $-2/\rho^{2}$, and the total-differential scale $s_1=-1/2$
contributes $s_1^{-1}=-2$; the net prefactor on every one of the eight words
is therefore
\begin{equation}\label{eq:ss:seednetprefactor}
\Big(-\frac{2}{\rho^{2}}\Big)(-2)\,\kappa_H^{2}\h_{\mathrm{HT}}\,
\frac{\rho^{2}}{\sqrt2}
=2\sqrt2\,\kappa_H^{2}\h_{\mathrm{HT}}
=\frac{\h g^{2}}{16\pi^{2}},
\end{equation}
with coefficient vector, in the ordered basis of~\eqref{eq:ss:seedunitresult},
\begin{equation}\label{eq:ss:seedmachinerow}
\Big(
\langle\bcW,\nabla_{+}\cW_{+}\rangle,\,
\langle\nabla_{+}\cW_{+},\bcW\rangle,\,
\langle\nabla_{+}\Phi_1,\widetilde\Phi^{1}\rangle,\,
\langle\widetilde\Phi^{1},\nabla_{+}\Phi_1\rangle,\,
\dots
\Big)=(1,-1,1,-1,1,-1,1,-1),
\end{equation}
matching~\eqref{eq:ss:seedunitresult} entry by entry. Equivalently, the
loop-scale relation
\begin{equation}\label{eq:ss:seedloopscale}
-2\,\kappa_H^{2}\h_{\mathrm{HT}}
=-\frac{\h g^{2}}{16\sqrt2\,\pi^{2}},
\qquad\text{i.e.}\qquad
\kappa_H^{2}\h_{\mathrm{HT}}
=\frac{\h g^{2}}{32\sqrt2\,\pi^{2}},
\end{equation}
identifies the twist computation with the boxed
result~\eqref{eq:ss:bbresult}: the matter words of~\eqref{eq:ss:seedbk},
which \cite{BK} obtains from the superpotential vertex, are the same matter
triangle as in paragraph~9. Only the zero-jet kernel entry
$K_{0,0;0,0}=1$ enters this channel; Budzik and Kulp's printed
master-triangle value is $\tfrac12$, exactly one half, the factor $2$ being
the Feynman-parameter prefactor of~\eqref{eq:ss:feynman}. The full tower
comparison is Section~\ref{sec:ht}.

\begin{remark}[Status of the comparison]\label{rem:ss:seedht}
The match of this paragraph is an external cross-check against~\cite{BK},
not a derivation input and not a completeness proof of the graph census: the
factor chain of paragraphs 1--10 was sealed before the twist formulas were
consulted. The loop-scale relation~\eqref{eq:ss:seedloopscale} uses the
conditional total-differential scale $s_1=-1/2$ of the dictionary. The
adversarial census review remains open; it does not alter the accepted
coefficient~\eqref{eq:ss:bbresult}.
\end{remark}
